\documentclass{article}
\usepackage[utf8]{inputenc}
\usepackage{amsmath}
\usepackage{graphicx}
\usepackage{accessibility}
\usepackage[numbers]{natbib}
\usepackage{amssymb}

\title{A Rust Implementation of Graph Parameter Algorithms for Parameterized Complexity Analysis}
\author{Jules}
\date{\today}

\begin{document}

\maketitle

\begin{abstract}
This paper details a Rust implementation of several graph parameter algorithms, inspired by the work of Komusiewicz et al. \cite{komusiewicz2025parameter}. We provide robust and efficient implementations for key parameters such as degeneracy and vertex cover, leveraging the \texttt{petgraph} library for graph data structures. This work serves as a practical, open-source companion to the theoretical analysis presented in the source paper, enabling further empirical studies in parameterized complexity.
\end{abstract}

\section{Introduction}
Parameterized algorithmics provides a framework for designing efficient algorithms for NP-hard problems by isolating the problem's hardness to a specific parameter \cite{downey1999parameterized}. The choice of parameter is crucial, and the work of Komusiewicz et al. \cite{komusiewicz2025parameter} provides a data-driven analysis of various graph parameters on real-world graphs. Motivated by their findings, this paper presents a Rust implementation of some of the key parameters discussed in their study. Our goal is to provide a high-quality, open-source library that can be used for further research and experimentation in this area.

\section{Background: Graph Parameters}
We implemented a selection of parameters from different categories as defined in Komusiewicz et al. \cite{komusiewicz2025parameter}.

\subsection{Degree-Related Parameters}
\paragraph{Degeneracy} The degeneracy of a graph is the smallest integer k such that every induced subgraph of the graph has a vertex of degree at most k. It can be computed in linear time.

\subsection{Modulator-Based Parameters}
\paragraph{Vertex Cover} A vertex cover is a subset of vertices such that each edge of the graph is incident to at least one vertex of the set. The vertex cover number is the size of the smallest such set. Finding the minimum vertex cover is NP-hard.

\subsection{Treewidth}
\paragraph{Treewidth} The treewidth of a graph is a measure of its tree-likeness. Computing it is NP-hard.

\section{Implementation in Rust}
Our implementation is a new module, \texttt{graph_theory}, within the \texttt{math_explorer} Rust library.

\subsection{Project Structure}
The module is organized as follows:
\begin{itemize}
    \item \texttt{graph.rs}: Defines the core graph data structure using \texttt{petgraph}.
    \item \texttt{parameters/degree.rs}: Implements degree-related parameters.
    \item \texttt{parameters/modulator.rs}: Implements modulator-based parameters.
    \item \texttt{parameters/treewidth.rs}: Implements treewidth.
\end{itemize}

\subsection{Data Structures}
We use the \texttt{petgraph} crate to provide the underlying graph data structures. Our \texttt{Graph} struct is a wrapper around \texttt{petgraph::Graph}, configured for undirected graphs.

\subsection{Algorithm Implementation}
\paragraph{Degeneracy} We implemented a linear-time algorithm to compute the degeneracy of a graph. The algorithm iteratively removes the vertex with the smallest degree and updates the degrees of its neighbors. The degeneracy is the maximum degree of a vertex at the time of its removal.

\paragraph{Vertex Cover} For vertex cover, we implemented a 2-approximation algorithm. This algorithm iterates through the edges of the graph. If an edge is not covered, it adds both of its endpoints to the cover. This is a well-known greedy algorithm that guarantees a solution at most twice the size of the optimal one \cite{bafna1999}.

\paragraph{Treewidth} The computation of treewidth is a notoriously difficult problem. The source paper \cite{komusiewicz2025parameter} uses specialized external solvers for this purpose. Our implementation provides a placeholder for this function, acknowledging its complexity and the need for external tools for an efficient solution.

\section{Conclusion}
This paper documents a Rust implementation of key graph parameters for parameterized complexity analysis. By providing a clean, well-tested, and open-source library, we hope to facilitate further research in this field. Future work could involve implementing more of the parameters from the source paper and integrating with external solvers for the NP-hard problems.

\bibliographystyle{plain}
\bibliography{references}

\end{document}
